\section*{Bash-Script}

Ist eine Textdatei die Befehle nacheinander ausführt und das Script so wieder verwendbar macht und Routine arbeiten automatisieren kann.

\begin{itemize}
  \item Es ist wie ein kleines Programm und wir bei \texttt{Unix}artigen Systemen verwendet
  \item Grundlegender Aufbau:
  \begin{itemize}
    \item \textbf{Shebang:} \texttt{\#!/bin/bash} teilt dem System mit welcher Interpreter verwendet werden soll zum ausführen des Codes.
    \item \textbf{Kommentare:} Sie werden verwendet mit \texttt{\#} um Code Informationen im Script für denn Menschen verständlich zu machen und wird vom System ignoriert.
    \item \textbf{Befehle:} Sind Systembefehle die nacheinander abgespielt werden und somit ein Abfolge von Aufgaben erledigen
  \end{itemize}
\end{itemize}

\begin{verbatim}
#!/bin/bash

# Dies ist ein Kommentar. Das Skript gibt eine Begrüßung aus.
echo "Hallo! Ich automatisiere jetzt deine Aufgabe."

# Eine Variable wird definiert und verwendet
BENUTZER=$(whoami)
echo "Du bist angemeldet als: $BENUTZER"

# Es wird geprüft, ob ein Ordner existiert, und ggf. angelegt
ORDNER="mein_backup_ordner"

if [ ! -d "$ORDNER" ]; then
    mkdir "$ORDNER"
    echo "Der Ordner '$ORDNER' wurde erfolgreich erstellt."
else
    echo "Der Ordner '$ORDNER' existiert bereits."
fi
\end{verbatim}

Es gibt genau so wichtige Bausteine wie in anderen Programmen bzw. Code wie \textbf{Variablen}, \textbf{Argumente}, \textbf{Bedingungen} und \textbf{Schleifen}.
